% 问题四
\section{问题四的模型建立与求解}

\subsection{问题分析}

问题四要求从优化蔬菜类商品补货和定价决策的角度出发，识别当前数据体系的不足，提出补充采集数据的建议，并论述这些数据对解决前述三个问题的具体帮助。

基于前三问的建模与求解过程，现有数据（附件1至附件4）存在以下五个方面的缺失，制约了各模型的进一步优化改进：

\begin{enumerate}
\item \textbf{日销量聚合口径无法区分共同趋势与真实关联。} 问题一中的Spearman相关系数基于日销量聚合计算，无法唯一识别购物篮层面的互补或替代关系。
\item \textbf{售价口径混合了正常销售与打折销售。} 附件2无法区分正常售价与傍晚清仓打折价，问题二中品类售价$P_{c,t}$为混合口径，导致价格弹性$\hat{\beta}_c$向零衰减。
\item \textbf{缺货日真实需求不可观测。} 附件2仅记录已售销量$Q_{i,t}$，问题三中基准需求$d_i$将系统性低估真实需求。
\item \textbf{损耗率为静态平均值，缺乏动态结构。} 附件4仅给出各商品近期的盘点周期平均损耗率$L_i$，问题二、三均假设损耗率恒定。
\item \textbf{数据粒度为日级别，缺少日内信息与外部参照。} 附件2为日粒度数据，无法支持日内动态调价；附件3缺少周边竞品零售价参照。
\end{enumerate}

基于上述诊断，本节从需求侧、市场侧、运营侧、供给侧与外部环境五个维度提出数据采集建议。


\subsection{数据采集建议}

\subsubsection{需求侧：真实会员ID与缺货登记日志}

\paragraph{采集内容}

在POS收银环节采集脱敏后的会员ID（或家庭户标识），建立跨期交易流水关联；同时，在收银台或线上小程序记录“消费者搜索但未找到”的缺货登记日志，记录单品$i$在日期$t$的缺货登记次数$m_{i,t}$。

\paragraph{模型帮助}

有了真实交易流水号后，可精准还原消费者的真实购物篮$\mathcal{B}_{m,t} = \{i : \text{会员 } m \text{ 在 } t \text{ 日购买了单品 } i\}$，直接开展购物篮关联分析，弥补问题一中仅能通过销量相关性刻画品类/单品关系的不足。计算单品$i$与$j$的提升度：
\begin{equation}
\text{Lift}(i \to j) = \frac{P(i \cap j)}{P(i)P(j)} = \frac{|\{m : i, j \in \mathcal{B}_m\}| / |\{m\}|}{\left(|\{m : i \in \mathcal{B}_m\}| / |\{m\}|\right) \left(|\{m : j \in \mathcal{B}_m\}| / |\{m\}|\right)}.
\end{equation}

当$\text{Lift}(i, j) > 1$时，$i$与$j$存在真实互补关系；$\text{Lift}(i, j) < 1$暗示替代关系。

利用会员ID可追踪消费者的跨期购买行为，将问题二、三中的“群体聚合幂律需求”升级为基于会员分层的需求模型（RFM分群+价格敏感度分层），实现分层差异化促销，而非“千人千面”个体定价（个人定价不现实且有合规风险）。

缺货登记数据可精确估计被抑制的真实需求。设真实需求满足：
\begin{equation}
D_{i,t}^* = Q_{i,t} + \alpha_i \cdot m_{i,t} + \varepsilon_{i,t},
\end{equation}
其中$\alpha_i$为单品$i$的缺货登记率，通过小规模抽样校准。直接修正问题三中单品基准需求$d_i$的估计偏差：
\begin{equation}
d_i^* = \frac{1}{T_{\text{win}}} \sum_{t \in \text{win}} \left[ Q_{i,t} + \alpha_i m_{i,t} \right],
\end{equation}
使报童模型中的$\mu_{c,t}$预测更为准确。

\paragraph{决策帮助}

补货量更贴近真实需求，降低缺货损失；按会员分层做差异化促销，兼顾利润与顾客留存；真实购物篮关联指导选品组合与关联陈列。

\paragraph{理由与可行性}

缺货登记近乎零成本（POS/小程序增加一个登记字段即可）；会员ID属中等改造成本，需做脱敏与隐私合规（《个人信息保护法》：最小化采集、匿名化、仅存聚合特征）。

\paragraph{偏误风险与验证}

风险：缺货登记可能重复/误报、会员数据稀疏。处理：按“会员+商品+时段”去重、设置最小登记量阈值；稀疏会员不做个体建模，仅聚合到分群。验证钩子：用缺货登记修正前后的$d_i$、$\mu_{c,t}$分别回测报童缺货率与利润，对比修正是否带来显著改善。


\subsubsection{市场侧：陈列拓扑图与竞品实时均价}

\paragraph{采集内容}

采集各蔬菜单品在货架上的物理位置坐标（空间距离矩阵），以及周边3公里内农贸市场、生鲜电商的同品类实时标价。

\paragraph{模型帮助}

问题三的需求函数$D_i(p)$假设单品需求彼此独立（仅通过品类覆盖约束软性关联），无法量化“小白菜涨价后消费者转向上海青”的替代效应。引入上述数据后，将问题三需求函数扩展为包含交叉价格弹性的形式：
\begin{equation}
D_i(p_i, \mathbf{p}_{-i}) = d_i \left( \frac{p_i}{p_i^0} \right)^{\beta_i} \prod_{j \ne i} \left( \frac{p_j}{p_j^0} \right)^{\epsilon_{ij}},
\end{equation}
其中$\epsilon_{ij}$为单品$i$对单品$j$价格的交叉弹性。$\epsilon_{ij}$通过历史价格变动事件估计，并依据货架距离$\text{dist}_{ij}$加权：
\begin{equation}
\epsilon_{ij} = \epsilon_{ij}^0 \cdot \exp(-\lambda \cdot \text{dist}_{ij}), \quad \lambda > 0.
\end{equation}
交叉弹性矩阵可直接嵌入问题三的MIP模型中，将目标函数中的$D_{i,k}$从单变量函数扩展为多变量耦合函数，使最优选品结果在替代品之间实现利润平衡，避免“同类相食（Cannibalization）”效应。当$\epsilon_{ij}=0$（$i\neq j$）时，式(4)自动退化为独立需求$D_i(p_i)=d_i(p_i/p_i^0)^{\beta_i}$，保证模型体系自洽。

引入竞品价格后，在定价约束中增设外部参照：
\begin{equation}
p_{i,k} \le (1 + \theta) \cdot P_{i,t}^{\text{comp}},
\end{equation}
其中$\theta$为商超相对竞品的可接受溢价率，由历史数据估计。

\paragraph{决策帮助}

定价考虑外部比价与内部替代，联合定价避免内部利润自相蚕食；陈列拓扑指导选品与货架布局，邻近替代品采用差异化定价。

\paragraph{理由与可行性}

现有附件无竞品数据，弹性估计存在竞争环境混淆；竞品价格可通过人工抽样或公开平台采集，成本中等，建议中期落地。

\paragraph{偏误风险与验证}

风险：竞品采样本底不一致、价格可比口径差异。处理：只对同口径价格建模、缺失用品类均值填充并打缺失标记。验证钩子：加竞品变量前后回归$R^2$与AIC对比；交叉弹性模型的选品利润与独立需求模型选品利润对比。


\subsubsection{运营侧：小时级库存水位与日内打折流水}

\paragraph{采集内容}

借助RFID或电子秤联网系统，采集日内分时段的库存余量与鲜度评级变化；同时记录黄昏清仓时段（如19:00后）的具体折扣率与对应销量（附件2已有“是否打折”标记，缺的是折扣率与时段，属低成本增量改造）；采集分时段客流量、客单价（可由收银小票聚合得到）。

\paragraph{模型帮助}

问题二、三假设每日为独立的“单周期报童问题”，未考虑日内价格随鲜度衰减的动态调整。引入上述时序数据后，可建立日内动态定价的随机控制模型。设每日分为$T$个时段，状态变量为$(I_\tau, g_\tau)$（$I_\tau$为剩余库存，$g_\tau$为品相等级），决策变量为折扣率$\delta_\tau$，值函数满足贝尔曼方程：
\begin{equation}
V_\tau(I_\tau, g_\tau) = \max_{\delta_\tau \in [0, \bar{\delta}]} \left\{ \mathbb{E}\left[ p(\delta_\tau)\min(D_\tau(p), I_\tau) + V_{\tau+1}\left(I_\tau - \min(D_\tau, I_\tau), g_{\tau+1}\right) \right] \right\},
\end{equation}
终端条件$V_T(\cdot, \cdot) = 0$。该模型确定各时段的最优折扣率，决定每日下午/傍晚的最优清仓时点与折扣梯度。

动态定价模型是对问题二、三中“固定日定价”策略的日内精细化补充（主模型仍是日定价，日内DP是补充层）。其输出的“预期日内折扣损失”可作为修正项，回传到报童模型的临界比$\kappa_{c,t}$中，避免因过度乐观定价导致晚间积压损耗。

\paragraph{决策帮助}

给出每日最优清仓时点与折扣梯度，平衡“晚间低价清货”与“利润损失”；补货量在临界比层面已预扣日内折价损失，减少期末报废。

\paragraph{理由与可行性}

打折流水为低成本增量改造（补折扣率与时段字段即可）；RFID/电子秤联网成本较高，作为长期可选项。

\paragraph{偏误风险与验证}

风险：时段库存测量噪声、鲜度评级主观、清仓折扣与正常售价混记。处理：鲜度评级用统一标准打分，折扣单单独标记。验证钩子：日内DP的期末损耗率与固定日定价回测对比；临界比修正前后的缺货率/报废率对比。


\subsubsection{供给侧：分环节损耗率与供应商履约指数}

\paragraph{采集内容}

按“运输 $\to$ 入库 $\to$ 上架 $\to$ 清仓”四个节点分段记录实际损耗率；同时追踪各供应商的历史到货提前期（LTD）及批次合格率；采集每日实时进货价与供应商报价（含波动区间）。

\paragraph{模型帮助}

问题二、三假设综合损耗率$L_i$在预测期内恒定（来自附件4）。引入分环节数据后，可将常数$L_i$拓展为基于历史分环节损耗率拟合的时变衰减函数。设单品$i$批次$v$在储存$\tau$小时后的环节$s$损耗率为：
\begin{equation}
L_{i,v}^{(s)}(\tau) = L_{i,0}^{(s)} \cdot \exp(\lambda_s \tau) \cdot \left(1 - \frac{\eta_{i,v}}{5}\right),
\end{equation}
其中$\lambda_s > 0$为环节$s$的损耗加速因子，由历史分环节损耗数据拟合得到，$\eta_{i,v}$为批次鲜度评分。综合损耗率为：
\begin{equation}
L_{i,v}(\tau) = 1 - \prod_{s=1}^{4} \left(1 - L_{i,v}^{(s)}(\tau)\right).
\end{equation}
时变损耗函数可替代问题二中的固定损耗放缩因子$1/(1-L_c)$，将补货量$R_{c,t}$的计算升级为：
\begin{equation}
R_{c,t} = \frac{\mu_{c,t}(P^*) + \Phi^{-1}(\kappa_{c,t}) \cdot \sigma_{c,t}(P^*)}{1 - \bar{L}_{c,t}},
\end{equation}
其中$\bar{L}_{c,t} = \sum_{i \in \mathcal{C}_c} w_i \cdot L_{i,v(t)}(\tau_t)$，$v(t)$为$t$日的供应批次编号，$\tau_t$为截止补货决策时该批次已储存的时长。

提前期不确定性引入安全库存项（需求率$\mu_d$、提前期标准差$\sigma_L$）：
\begin{equation}
\text{SS} = z \cdot \sqrt{\sigma_d^2 \cdot L + \mu_d^2 \cdot \sigma_L^2}.
\end{equation}
实时批发价可将“未来一周批发价取最近均价”的假设升级为价格情景区间，支撑灵敏度分析。

\paragraph{决策帮助}

补货量同时反映“到货品质批次差异”和“上架后随时间衰减”的双重不确定性，显著提升报童模型在恶劣供应链条件下的抗风险能力；供应商履约指数可指导供应商选择与订货量分配。

\paragraph{理由与可行性}

分环节损耗需门店交接单改造，成本中低，建议短期优先；供应商履约数据需与供货方对接，属中期项。

\paragraph{偏误风险与验证}

风险：环节归属不清、提前期样本少。处理：交接单逐节点签字确认损耗责任；提前期至少收集10个批次再启用安全库存公式。验证钩子：时变损耗与恒定损耗的期末损耗率回测对比；含/不含提前期安全库存的缺货率对比。


\subsubsection{外部环境：天气、节假日与促销日历}

\paragraph{采集内容}

接入日/小时级天气数据（温度、降水量、湿度、极端天气哑变量）；整理节假日、节气、周末与大型活动日历；记录商超自身促销活动计划（促销品类、折扣力度、起止时间）。

\paragraph{模型帮助}

天气与节假日是蔬菜需求最强的外生驱动因素（4–10月供给丰富、极端天气推高叶菜需求），直接进入问题二的需求—价格回归与Prophet预测，作为外生回归项：
\begin{equation}
\begin{split}
\ln Q_{c,t} = \;&\alpha_c + \beta_c\ln P_{c,t} \\
&+ \sum_{k=1}^{6}\delta_k\mathbf{1}_{\{\text{weekday}=k\}}
+ \sum_{m=1}^{11}\eta_m\mathbf{1}_{\{\text{month}=m\}}
+ \gamma_T T_{c,t} + \gamma_R R_{c,t} \\
&+ \theta_1 \text{Holiday}_t + \theta_2 \text{Promo}_t + \varepsilon_{c,t},
\end{split}
\end{equation}
其中$T_{c,t}$为温度，$R_{c,t}$为降水量，$\text{Holiday}_t$为节假日哑变量，$\text{Promo}_t$为促销哑变量。

引入这些变量后，可剔除天气和促销活动对销量的干扰，使$\hat{\beta}_c$的估计更加准确；同时提升Prophet预测中基准需求$\bar{Q}_{c,t}$与残差标准差$\sigma_{c,t}$的精度，改善报童补货量的均值—方差结构。

\paragraph{决策帮助}

预测期若遇极端天气或节假日，自动上调补货量与价格保护区间；促销日历与定价联动，避免促销期与非促销期需求混淆导致弹性误估。

\paragraph{理由与可行性}

成本极低（公开API+内部日历），见效最快，建议最高优先级。

\paragraph{偏误风险与验证}

风险：天气数据粒度不一致、促销与需求互为因果。处理：用门店最近气象站数据；促销变量只用于“去混淆”，弹性估计以非促销样本为主或加入促销交互项。验证钩子：加天气/促销变量前后Prophet回测RMSE与MAE对比。


\subsection{综合分析}

\subsubsection{数据—模型—决策对应总览}

\begin{table}[htbp]
\centering
\caption{数据—模型—决策对应关系}
\label{tab:mapping}
\begin{tabular}{cccc}
\toprule
数据类别 & 影响的前问模型/参数 & 改善的决策 & 实施成本 \\
\midrule
会员ID/购物篮 & Q1 关联分析 & 选品组合、关联陈列 & 中 \\
缺货登记 & Q2 $\mu_{c,t}$、Q3 $d_i$ & 补货量、缺货率 & 低 \\
会员分层 & Q2 需求分层、Q3 单品需求 & 分群促销 & 中 \\
竞品实时均价 & Q2 $\beta_c$校准、Q3 $\epsilon_{ij}$ & 联合定价 & 中 \\
陈列拓扑 & Q3 替代/互补结构 & 货架布局、选品 & 中 \\
日内库存/鲜度 & Q2/Q3 $\kappa_{c,t}$、日内DP & 清仓时点、折扣梯度 & 高 \\
日内打折流水 & Q2/Q3 $\kappa_{c,t}$ & 清仓折扣、期末报废 & 低 \\
分时段客流 & Q2 日内需求估计 & 排班、促销时段 & 中高 \\
分环节损耗率 & Q2/Q3 $L_i \to \bar{L}_{c,t}$ & 补货量、报废率 & 中低 \\
供应商履约 & Q2 安全库存 & 供应商选择 & 中 \\
实时批发价 & Q2 $W_{c,t}$情景化 & 定价调整 & 中低 \\
天气/节假日 & Q2 回归/Prophet外生项 & 补货与价格预警 & 极低 \\
促销日历 & Q2 弹性去混淆 & 促销与定价联动 & 低 \\
\bottomrule
\end{tabular}
\end{table}

\subsubsection{实施优先级与ROI评估}

五类数据在采集成本和模型改进效果上存在差异。采集价值可量化为“决策增益 − 采集成本”：P1类数据（天气、缺货登记、促销日历、分环节损耗）总改造成本最低，而预期缺货/报废率下降带来的毛利提升通常远高于成本，短期投入回报为正。优先级排序见表\ref{tab:priority}。

\begin{table}[htbp]
\centering
\caption{五类数据采集方案的优先级排序}
\label{tab:priority}
\begin{tabular}{ccccc}
\toprule
优先级 & 数据维度 & 数据方案 & 采集方式 & 实施成本 \\
\midrule
1 & 外部环境 & 天气、节假日与促销日历 & API接入+内部排期表 & 极低 \\
2 & 需求侧 & 会员ID+缺货登记 & POS系统改造+收银台登记 & 中 \\
3 & 运营侧 & 折扣标记+分时段库存 & POS/电子秤系统升级 & 低 \\
4 & 供给侧 & 分环节损耗率+供应商评级 & 验收流程+PDA录入 & 中 \\
5 & 市场侧 & 货架位置+竞品均价 & 平面测绘+数据采集 & 高 \\
\bottomrule
\end{tabular}
\end{table}

\subsubsection{数据间的协同效应}

五类数据并非独立发挥作用：

\begin{enumerate}
\item 天气/促销日历与折扣标记结合：天气和促销日历可作为外生变量辅助区分“促销导致的销量变化”与“正常需求波动”，使折扣标记$r_{i,t}$的分层弹性估计更加准确；
\item 折扣标记与缺货登记结合：折扣标记$r_{i,t}$刻画价格降低对销量的影响，缺货登记$m_{i,t}$刻画价格过低导致缺货时被抑制的需求，二者共同覆盖需求函数$D_i(p)$的可观测区间两端；
\item 购物篮关联与货架距离结合：购物篮关联$\text{Lift}(i,j)$与货架距离$\text{dist}_{ij}$结合，可估计空间距离对替代弹性的衰减因子$\lambda$；
\item 分时段库存与分环节损耗率结合：分时段库存$I_{i,t,\tau}$与分环节损耗率$L_{i,v}^{(s)}(\tau)$结合，损耗率随时间加速增长的函数关系可直接嵌入贝尔曼方程的状态转移。
\end{enumerate}

\subsection{小结}

问题四从需求侧、市场侧、运营侧、供给侧与外部环境五个维度提出了五类补充数据的采集方案，分别对应前三个模型的关键局限：

\begin{enumerate}
\item 需求侧——会员ID与缺货登记：解决问题一销量相关无法识别真实购物篮关联的问题，修正问题三因缺货截断导致的需求低估；
\item 市场侧——陈列拓扑与竞品均价：为问题三引入交叉价格弹性和外部竞争定价约束，弥补独立需求假设的不足；
\item 运营侧——日内库存与打折流水：解决问题二价格弹性估计中混合口径导致的衰减偏误，并将日清报童扩展为日内动态定价；
\item 供给侧——分环节损耗率与供应商评级：将问题二、三的恒定损耗率替换为基于批次品质和储存时长的动态损耗函数；
\item 外部环境——天气、节假日与促销日历：提升需求预测精度，剔除促销和天气对弹性估计的干扰。
\end{enumerate}